\subsection{Setup and Formal Definition of the Tokenisation Operator}

To rigorously understand the failure modes of patch-based architectures like Chronos-Bolt, we must first formally define the tokenisation process. Let $x[n]$ represent a discrete-time, univariate signal where the index $n \in \mathbb{Z}$ is sampled at a continuous frequency of $f_s$.

Unlike classical autoregressive models that process temporal data one scalar time-step at a time, modern foundation models employ patch-based tokenisation. This technique groups contiguous samples to radically reduce the effective sequence length and improve the computational efficiency of the transformer's attention mechanisms. This is achieved by sliding a window of fixed size $P$ (the patch length) across the signal, moving forward by a step size of $S$ (the stride).

We can extract this sequence of tokens as follows:
\begin{equation}
    \mathbf{p}_k = \bigl(x[kS],\; x[kS+1],\; \ldots,\; x[kS+P-1]\bigr) \in \mathbb{R}^P, \quad k \in \mathbb{N}
\end{equation}

\begin{definition}[Patch Tokenisation Operator]
    We can formalise this extraction as a linear map acting on the bounded signal space. The map $\mathcal{T}: \ell^\infty(\mathbb{Z}) \to (\mathbb{R}^P)^{\mathbb{N}}$, as defined above, is referred to as the \emph{patch tokenisation operator}. When collecting an arbitrary number of $K$ patches from a signal vector $\mathbf{x}$, the operation can be represented as a matrix multiplication:
    \begin{equation}
        \mathbf{P} = \mathcal{T}\{x\} = \mathbf{W} \cdot \mathbf{x}
    \end{equation}
    Here, $\mathbf{W} \in \{0,1\}^{KP \times N}$ is a block-binary selection matrix. Its $k$-th block consists of $P$ rows that act as a mask, picking out the exactly $P$ consecutive signal samples starting at the temporal index $kS$.
\end{definition}

\subsection{The Mechanics of Phase-Locking and Degeneracy}

The fundamental vulnerability of the patch tokenisation operator emerges when the periodic structure of the input signal aligns perfectly with the model's fixed stride grid. We refer to this specific synchronization as ``phase-locking.''

\begin{proposition}[Condition for Token Identity]
    Two consecutive patches, $\mathbf{p}_k$ and $\mathbf{p}_{k+1}$, will be mathematically identical if and only if the signal amplitude remains constant across the stride shift for every element within the observation window:
    \begin{equation}
        x[n] = x[n + S] \quad \forall n \in \{kS,\, kS+1,\, \ldots,\, kS+P-1\}
    \end{equation}
    Globally, if the underlying continuous signal is a pure sinusoid (or any periodic signal) with a fundamental period of $T_0$ samples, this exact alignment occurs whenever the stride $S$ spans an integer multiple of the signal's period:
    \begin{equation}
        \boxed{S = m \cdot T_0, \quad m \in \mathbb{Z}^+}
    \end{equation}
    By substituting $T_0 = f_s / f$, we can express this condition in terms of the signal's continuous frequency $f$ and the sampling frequency $f_s$:
    \begin{equation}
        f = m \cdot \frac{f_s}{S}, \quad m \in \mathbb{Z}^+
    \end{equation}
\end{proposition}

\begin{proof}
    To satisfy $\mathbf{p}_k = \mathbf{p}_{k+1}$ strictly component-wise, it is required that $x[kS + j] = x[(k+1)S + j]$ for every internal patch index $j \in \{0,\ldots,P-1\}$. For a periodic signal defined by $T_0 = f_s/f$, this strict equality holds true whenever the stride shift $S$ represents a traversal of exactly $m$ full wave cycles. In such instances, the sliding window captures the exact same geometrical segment of the wave at every step.
\end{proof}

\subsection{The Consequence: Loss of Observability}

When the phase-locking condition $f = m \cdot (f_s/S)$ is met, we achieve a state where $\mathbf{p}_k = \mathbf{p}_{k+1}$ for all possible values of $k$. The token matrix $\mathbf{P}$ undergoes a mathematical rank collapse, reducing its rank to exactly 1.

This creates a ``degenerate token sequence'', which induces severe representational failures within the network:

\begin{itemize}
    \item \textbf{Complete Temporal Blindness:} Because the transformer processes a sequence of completely identical vectors, the temporal variation of the actual underlying signal becomes entirely invisible to the model's self-attention mechanisms. The model cannot perceive that a wave is oscillating; it simply perceives a constant, flat embedding.
    \item \textbf{Observational Equivalence and Aliasing:} Any two distinct signals that differ only by a phase shift of $S$ samples will produce the exact same sequence of tokens. They become observationally indistinguishable to the decoder, making it impossible for the model to correctly reconstruct or forecast the true phase of the original signal.
    \item \textbf{Soft Degradation at Boundary Frequencies:} This failure is not isolated to exact integer multiples. Even when patches are not perfectly identical but are merely highly correlated (i.e., when $S \approx m T_0$), the sequence suffers a \emph{soft} loss of observability. The effective information bandwidth is severely restricted, explaining why predictive accuracy degrades smoothly around these critical frequency boundaries before fully collapsing.
\end{itemize}

For example, in the specific architecture of Chronos-Bolt, the sampling frequency is $f_s = 512$ Hz and the stride is $S=16$. According to our derived condition $f = m \cdot (512/16)$, the fundamental patch frequency is precisely 32 Hz. Consequently, any input signal at 32 Hz, 64 Hz, 96 Hz, or other integer harmonics will trigger this precise structural aliasing, resulting in an immediate and sharp collapse in the model's predictive accuracy.